% Copyright (c) 2022-2026 MKM Research Labs. Licensed under the MIT License; see LICENSE.
%
% MKM-FPO-001 Flood Polynomial Model -- model documentation v2.0, the document of record
% filed on the model as Core Documentation. It supersedes v1.0 (28-Mar-2026). The Wolfram
% benchmark of v1.0 (MKM-FPO-001_wolfram_benchmark v0.1) found the model fell as the water
% rose and that none of the document's printed tables or accuracy claims held; v2.0 is a
% new model, refitted, and this document describes it. Its tables are inlined, generated
% from the code and the fit.
%
% Build: it inputs the house header, which lives in PhysicalRisk. Copy
% PhysicalRisk/docs/models/shared beside this folder as ../shared, then run
% pdflatex three times (the third settles "Page n of N").
\documentclass[11pt,a4paper]{article}

\newcommand{\doctitle}{Flood Polynomial Model}
\newcommand{\docsubtitle}{MKM-FPO-001 --- A Closed-Form Surrogate for the Gauge Flood Classifiers}
\newcommand{\docversion}{2.0}
\newcommand{\docdate}{03-November-2026}
\newcommand{\docauthor}{David K Kelly}

\input{../shared/mkm_header}
% T1 so the underscores in code identifiers print as underscores.
\usepackage[T1]{fontenc}
\usepackage{lmodern}
% The shared history table sets every column to its natural width; wrap the description.
\renewcommand{\mkmhistorypage}[1]{%
    \section{Document History}
    \begin{table}[ht]
        \centering
        \begin{tabular}{l|l|p{9cm}|l}
            \textbf{Date} & \textbf{Version} & \textbf{Description} & \textbf{Author} \\
            \hline
            #1
        \end{tabular}
    \end{table}
}

\newcommand{\fpNearMaeLo}{0.0408}
\newcommand{\fpNearMaeHi}{0.0658}
\newcommand{\fpNearMaeMean}{0.0544}
\newcommand{\fpNearBiasLo}{$-$0.034}
\newcommand{\fpNearBiasHi}{0.025}
\newcommand{\fpNearOverLo}{10}
\newcommand{\fpNearOverHi}{21}
\newcommand{\fpBelowMaeLo}{0.0064}
\newcommand{\fpBelowMaeHi}{0.0203}
\newcommand{\fpFitStates}{1{,}807{,}233}
\newcommand{\fpTestStates}{1{,}807{,}896}
\newcommand{\fpMinSlope}{0.0100}

\begin{document}
% Long code paths and identifiers cannot hyphenate; let the spacing stretch instead.
\sloppy

\mkmtitlepage
\mkmlegalpage

\tableofcontents
\clearpage

% ===================================================================
\mkmexecsummary
% ===================================================================

The Flood Polynomial Model (MKM-FPO-001) gives the probability that a gauge goes on to
flood, from its water level and the hour of the storm, in closed form. It stands in for
the per-gauge gradient-boosted classifiers (MKM-FC-001) where a stress scenario needs
$P(\text{flood})$ without loading a trained model for every gauge.

This is version 2.0, a new model. Version 1.0 was benchmarked against an independent
Wolfram calculation and against the classifiers it approximates, and it failed. At
hour~24 it fell as the water rose between 52\% and 74\% of the severe level. At the
severe level it gave 0.00 at hour~0 and 0.29 at hour~167, where its document printed
0.42 and 0.97. Eight of the 54 cells of its printed table were reproduced by the code,
and none of its accuracy figures held. It had been fitted to one gauge.

Version 2.0 changes the model in four ways:
\begin{itemize}[nosep]
    \item \textbf{At or above severe, $P=1$.} Reaching the severe level is the flood
          the classifiers are trained to predict. Below it, $P$ is a fitted logistic.
    \item \textbf{Fitted to every gauge.} One kernel, cubic in the two features,
          fitted to all 11 gauges' classifiers on \fpFitStates{} real storm states and
          measured on \fpTestStates{} held-out ones.
    \item \textbf{Rising water never lowers $P$.} The fit is constrained to make $P$
          non-decreasing in the level at every hour.
    \item \textbf{No extrapolation early in a storm.} No storm reaches 40\% of severe
          before hour~4.7, or severe before hour~37.4. Earlier than that the
          classifiers have never seen a state, so the model holds the hour at that
          envelope instead of extrapolating.
\end{itemize}

On held-out storms near severe ($w/s\ge0.75$), the mean absolute difference from each
gauge's classifier is \fpNearMaeLo{} to \fpNearMaeHi{} across the 11 gauges, and the
bias is \fpNearBiasLo{} to \fpNearBiasHi{}. The difference exceeds 0.1 in
\fpNearOverLo\% to \fpNearOverHi\% of those states. That is the price of a formula that
sees only the level and the hour: the classifiers also see how fast the water is
rising, and no function of level and hour can follow that.

% ===================================================================
\mkmpurpose
% ===================================================================

\subsection{Purpose}
The model answers one question in closed form: given a gauge at water level $w$ in
hour $n$ of a storm, how likely is it to reach its severe threshold $s$? The stress
scenarios ask it hour by hour along a hydrograph, for many gauges, where loading a
trained classifier per gauge is not practical.

\subsection{Scope}
\begin{itemize}[nosep]
    \item \textbf{In scope:} $P(\text{flood})$ from $w/s$ and the storm hour, for any
          gauge with a severe threshold; its slope in the level, for sensitivity.
    \item \textbf{Out of scope:} the rate of rise and other features the classifiers
          use; gauge-specific behaviour; spatial correlation between gauges.
\end{itemize}

\subsection{How Production Uses It}
\begin{itemize}[nosep]
    \item \textbf{Single-gauge stress scenario}
          (\texttt{src/routes/trading/stress/scenario.py}) evaluates
          \texttt{p\_flood\_poly} at each hour of the scenario hydrograph, and prints a
          probability surface over levels and hours. Once the level has reached
          severe, the scenario holds $P=1$ for the rest of the storm, even if the water
          falls back: the flood has happened, and the swaps it triggers are knocked
          out. That rule is the scenario's, not the model's.
    \item \textbf{Portfolio stress screen}
          (\texttt{src/static/js/trading/port\_stress/setup.js}) draws $P(\text{flood})$
          hour by hour in the browser. It fetches the model from
          \texttt{GET /api/v1/trading/stress/flood-poly} and evaluates it there. Until
          v2.0 it carried its own copy of v1.0's coefficients.
    \item \textbf{Portfolio stress P\&L}
          (\texttt{src/routes/trading/port\_stress/\_compute.py}) uses
          \texttt{p\_flood\_simple}, which is $\min(1, w/s)$. That is a proximity ratio,
          not a probability (\S\ref{sec:simple}).
    \item \textbf{Sensitivity} (\texttt{src/models/stress/sensitivity.py}) evaluates
          the model and its slope over candidate errors in the severe threshold.
\end{itemize}

\subsection{Model Chain Position}
\[
\text{MKM-FC-001 (per-gauge classifiers)}
\xrightarrow{\text{fitted to}}
\textbf{MKM-FPO-001}
\to \text{stress scenarios (MKM-ST-001)}
\]

% ===================================================================
\mkmframework
% ===================================================================

\subsection{Features}
Two gauge-independent features, as the classifiers use them:
\begin{align}
h &= \ln\!\left(\max\!\left(\frac{w}{s},\; \phi\right)\right), \qquad \phi = 0.2,
    \label{eq:h} \\
t &= \ln\!\left(\frac{n+1}{T}\right), \qquad T = 168, \quad n\in\{0,\dots,167\}.
    \label{eq:t}
\end{align}
$h$ is zero at the severe level and negative below it. $t$ runs from
$\ln(1/168)\approx-5.12$ at hour~0 to zero at hour~167. Below the level floor
$\phi$, 20\% of severe, the level is held at the floor.

\subsection{Model Equation}
\begin{equation}
\label{eq:model}
P(\text{flood}) =
\begin{cases}
1 & w \ge s, \\[2pt]
\sigma\big(z(h,\ \max(t,\ \tau(h)))\big) & w < s,
\end{cases}
\qquad
z(h,t) = \sum_{k=1}^{10} c_k\, h^{i_k} t^{j_k},
\end{equation}
where $\sigma(z) = 1/(1+e^{-z})$ with $z$ clamped to $[-30, 30]$, the ten terms
$(i_k, j_k)$ are every power of $h$ and $t$ up to total degree three, and
$\tau(h)$ is the support envelope of \S\ref{sec:envelope}.

\textbf{At or above severe.} A gauge at its severe level has flooded: that is the
event the classifiers are trained to predict. So $P=1$ there, at every hour. Version
1.0 had no such rule and fell back to 0.29 at the severe level late in a storm; the
stress route had been setting $P=1$ there itself, without the document saying so.

\textbf{Below severe.} The cubic kernel is fitted to the classifiers
(\S\ref{sec:calibration}). From hour~16 on, at a fixed level, $P$ never rises as the
storm goes on: a level reached late in a storm is one the water has been slow to reach,
and it is less likely to go on to severe. That is what the classifiers say, and the
reverse of what v1.0's document described. In the first 16 hours, just after the
envelope, $P$ can rise with the hour before it falls, by at most 0.06.

\subsection{The Support Envelope}\label{sec:envelope}
The kernel is fitted where storms go. Storms rise at a finite rate, so a level is
first reached some hours into a storm, and the classifiers have never been shown a
state earlier than that. There, the polynomial is extrapolating, and a cubic
extrapolates freely: refitted without a guard, it gave 0.998 at half of severe at
hour~0, where no storm has ever been.

$\tau(h)$ is the earliest log time at which any gauge's storms reached the level $h$,
measured over 2,000 storms at each of the 11 gauges (Table~\ref{tab:envelope}), and
interpolated linearly in $t$ between knots. Below the first knot, 32.6\% of severe,
every hour is supported. Earlier than the envelope, the model holds the hour there:
a level reached sooner than any storm reached it is priced as the earliest storm that
did. The cells this affects are marked in Table~\ref{tab:grid}.

\begin{table}[ht]\centering\begin{tabular}{rr|rr}
$w/s$ & earliest hour & $w/s$ & earliest hour \\ \hline
0.326 & 0.0 & 0.7 & 15.1 \\
0.35 & 2.3 & 0.75 & 18.4 \\
0.4 & 4.7 & 0.8 & 22.0 \\
0.45 & 6.1 & 0.85 & 24.9 \\
0.5 & 7.5 & 0.9 & 29.3 \\
0.55 & 8.9 & 0.95 & 35.7 \\
0.6 & 10.8 & 1 & 37.4 \\
0.65 & 12.9 &  &  \\
\end{tabular}\caption{The support envelope: the earliest storm hour at which any gauge's storms reached each fraction of severe. Earlier than this the hour is held here; between knots, linearly in $t$.}\label{tab:envelope}\end{table}

\subsection{Monotonicity}
Rising water never lowers $P$, at any hour. Three parts make that hold:
\begin{itemize}[nosep]
    \item Between the floor and severe, the fit holds $\partial z/\partial h \ge 0.01$
          on a grid of 120 levels and all 168 hours.
    \item Where the hour is held at the envelope, $P$ moves with the level along the
          envelope, at rate
          $\partial z/\partial h + \partial z/\partial t \cdot \tau'(h)$. The fit holds
          that at least 0.01 too, at 20 points on each segment.
    \item Below the floor $P$ is constant, and at severe it steps up to 1.
\end{itemize}
The smallest constrained slope in the fit is \fpMinSlope{}. The test suite checks $P$
is non-decreasing over 2,000 levels at every one of the 168 hours.

The hour is a whole number (Eq.~\eqref{eq:t}), and every caller passes one. That
matters in one place. Between hour~0 and hour~1, where $t$ moves fastest, the kernel's
slope in the level dips below zero near a third of severe (to about $-0.10$ at
hour~0.4). No whole hour falls there, so the model never evaluates it, but the kernel
is not monotone at fractional hours and should not be called with one.

% ===================================================================
\mkminputs
% ===================================================================

\subsection{Calibration Data}\label{sec:calibration}
The target is each gauge's production classifier, read as production reads it, on the
real storm states it is trained on: 2,000 storms per gauge, drawn once with seed~7 by
the challenger's catalogue loader, for all 11 gauges with a classifier. Only states
below severe are used, because at and above it $P=1$ by definition. Half of each
gauge's storms are fitted and the other half measured. That gives \fpFitStates{}
fitted and \fpTestStates{} held-out states.

\subsection{Fitting}
\texttt{python -m tools.flood\_poly\_fit} minimises the weighted squared difference
between $\sigma(z)$ and the classifier's probability, pooled over every gauge, with
the monotonicity penalty above. States at or above 75\% of severe weigh five times as
much, because that is where a stress scenario's outcome is decided. The fitted states
are evaluated as the model evaluates them, with the level floor and the envelope
applied.

The optimiser's tolerances matter. At SciPy's defaults ($10^{-8}$) the fit stopped on a
plateau after 94 evaluations, and on synthetic data it failed to recover a kernel it
had been shown. At $10^{-15}$ it recovers it exactly. The tolerances are in
\texttt{config.models} and the test suite checks the recovery.

\subsection{Coefficients}
\begin{table}[ht]\centering\begin{tabular}{c|c|r}
$k$ & Term $h^i t^j$ & $c_k$ \\ \hline
1 & $h$ & $+32.871867$ \\
2 & $t$ & $-12.801048$ \\
3 & $h t$ & $-7.540193$ \\
4 & $h^2$ & $+44.826451$ \\
5 & $t^2$ & $-2.848226$ \\
6 & $h^3$ & $+54.512653$ \\
7 & $t^3$ & $-0.756589$ \\
8 & $h^2 t$ & $-29.568283$ \\
9 & $h t^2$ & $+6.254275$ \\
10 & $1$ & $-1.293810$ \\
\end{tabular}\caption{Fitted coefficients, as \texttt{config.models} holds them.}\label{tab:coeffs}\end{table}

\subsection{Model Parameters}
\begin{center}
\begin{tabular}{llr}
\toprule
\textbf{Parameter} & \textbf{Config name} & \textbf{Value} \\
\midrule
Level floor $\phi$ & \texttt{flood\_poly\_level\_floor} & 0.2 \\
Storm horizon $T$ & \texttt{flood\_poly\_storm\_hours} & 168 \\
Exponent clamp & \texttt{flood\_poly\_exp\_clamp} & 30 \\
Kernel terms & \texttt{flood\_poly\_terms} & 10 \\
Support envelope & \texttt{flood\_poly\_envelope} & 15 knots \\
\bottomrule
\end{tabular}
\end{center}

% ===================================================================
\section{Fit Quality}
% ===================================================================

Measured on the held-out storms, against each gauge's own classifier:

\begin{table}[ht]\centering\begin{tabular}{l|rrr|r}
Gauge & Near MAE & Near bias & Near $>0.1$ & Below MAE \\ \hline
\texttt{GAUGE-085522ee} & 0.0590 & $+0.0247$ & 20\% & 0.0125 \\
\texttt{GAUGE-1aaf8b08} & 0.0408 & $-0.0070$ & 10\% & 0.0087 \\
\texttt{GAUGE-9e5f3ef9} & 0.0542 & $+0.0128$ & 16\% & 0.0198 \\
\texttt{GAUGE-9fffb377} & 0.0658 & $-0.0343$ & 21\% & 0.0203 \\
\texttt{GAUGE-a6be4ee9} & 0.0465 & $+0.0171$ & 15\% & 0.0075 \\
\texttt{GAUGE-b2dbe58b} & 0.0518 & $+0.0087$ & 17\% & 0.0100 \\
\texttt{GAUGE-c0fbc443} & 0.0478 & $-0.0015$ & 15\% & 0.0064 \\
\texttt{GAUGE-d208bf3d} & 0.0557 & $-0.0081$ & 16\% & 0.0179 \\
\texttt{GAUGE-dc5e88fb} & 0.0578 & $+0.0246$ & 21\% & 0.0092 \\
\texttt{GAUGE-ef53398c} & 0.0628 & $-0.0288$ & 20\% & 0.0166 \\
\texttt{GAUGE-f8629444} & 0.0566 & $+0.0188$ & 18\% & 0.0174 \\
\end{tabular}\caption{Agreement with each gauge's GBM on the held-out storms (\fpTestStates{} states). ``Near'' is $w/s\ge 0.75$; ``$>0.1$'' is the share of near states where $|P-\mathrm{GBM}|$ exceeds 0.1.}\label{tab:fit}\end{table}

Below severe overall the mean absolute difference is \fpBelowMaeLo{} to
\fpBelowMaeHi{}, because most states are far below severe and both models give
nearly zero there. Near severe the difference is \fpNearMaeLo{} to \fpNearMaeHi{}. The
bias near severe is \fpNearBiasLo{} to \fpNearBiasHi{}: the model is not conservative
everywhere, and it under-predicts at some gauges.

The difference that remains is mostly the rate of rise. The classifiers give
different probabilities to two states at the same level and hour, one rising fast and
one barely rising. A function of level and hour can at best give their average.

% ===================================================================
\mkmimplementation
% ===================================================================

\subsection{Source Code}
\begin{itemize}[nosep]
    \item \texttt{src/models/stress/flood\_poly.py}: \texttt{p\_flood\_poly},
          \texttt{p\_flood\_series}, \texttt{p\_flood\_simple}, the features and
          envelope, \texttt{model\_card}, \texttt{client\_params}.
    \item \texttt{config/models/\_misc.py}: coefficients, terms, floor, envelope and
          fit settings.
    \item \texttt{src/routes/trading/stress/flood\_poly.py}: serves
          \texttt{client\_params} to the browser.
    \item \texttt{src/static/js/trading/port\_stress/setup.js}: the browser's
          evaluation of the served model.
    \item \texttt{tools/flood\_poly\_fit/}: the fit.
\end{itemize}

\subsection{Usage}
\begin{lstlisting}[language=Python]
from models.stress.flood_poly import p_flood_poly, p_flood_series

p = p_flood_poly(water_level=5.2, hour=48, severe_level=6.58)
ps = p_flood_series([4.0, 4.5, 5.2, 5.8, 5.5], severe_level=6.58)  # hour = position
\end{lstlisting}

\subsection{\texttt{p\_flood\_simple} Is Not a Probability}\label{sec:simple}
\texttt{p\_flood\_simple(w, s)} returns $\min(1, w/s)$. It is a proximity measure:
a gauge at half its severe level scores 0.5, where Table~\ref{tab:grid} gives a
probability anywhere from nearly zero to 0.66 depending on the hour. Portfolio stress
P\&L uses it where only a gauge's peak level is known. Its name and v1.0's document
presented it as a time-independent $P(\text{flood})$; it is kept, under that name, as
a ratio, and is documented as one.

\subsection{Numerical Safeguards}
\begin{itemize}[nosep]
    \item The level floor keeps $\ln(w/s)$ finite at zero or negative water levels.
    \item The exponent is clamped to $[-30, 30]$ before the logistic.
    \item A severe threshold at or below zero gives $P = 0$.
\end{itemize}

\subsection{Dependencies}
The model uses only the Python standard library at run time.

% ===================================================================
\mkmsensitivity
% ===================================================================

\subsection{Probability Surface}
\begin{table}[ht]\centering\small\begin{tabular}{r|rrrrrrrr}
$w/s$ & hr 0 & hr 6 & hr 12 & hr 24 & hr 48 & hr 84 & hr 120 & hr 167 \\ \hline
0.20 & 0.001 & 0.000 & 0.000 & 0.000 & 0.000 & 0.000 & 0.000 & 0.000 \\
0.30 & 0.046 & 0.029 & 0.001 & 0.000 & 0.000 & 0.000 & 0.000 & 0.000 \\
0.40 & 0.334$^\dagger$ & 0.388 & 0.352 & 0.043 & 0.000 & 0.000 & 0.000 & 0.000 \\
0.50 & 0.608$^\dagger$ & 0.608$^\dagger$ & 0.664 & 0.490 & 0.034 & 0.000 & 0.000 & 0.000 \\
0.60 & 0.896$^\dagger$ & 0.896$^\dagger$ & 0.891 & 0.801 & 0.293 & 0.008 & 0.000 & 0.000 \\
0.70 & 0.990$^\dagger$ & 0.990$^\dagger$ & 0.990$^\dagger$ & 0.969 & 0.733 & 0.079 & 0.003 & 0.000 \\
0.80 & 0.999$^\dagger$ & 0.999$^\dagger$ & 0.999$^\dagger$ & 0.999 & 0.970 & 0.461 & 0.036 & 0.001 \\
0.90 & 1.000$^\dagger$ & 1.000$^\dagger$ & 1.000$^\dagger$ & 1.000$^\dagger$ & 0.999 & 0.937 & 0.348 & 0.013 \\
0.95 & 1.000$^\dagger$ & 1.000$^\dagger$ & 1.000$^\dagger$ & 1.000$^\dagger$ & 1.000 & 0.988 & 0.714 & 0.054 \\
0.99 & 1.000$^\dagger$ & 1.000$^\dagger$ & 1.000$^\dagger$ & 1.000$^\dagger$ & 1.000 & 0.997 & 0.906 & 0.165 \\
1.00 & 1.000 & 1.000 & 1.000 & 1.000 & 1.000 & 1.000 & 1.000 & 1.000 \\
\end{tabular}\caption{$P(\mathrm{flood})$ by level and storm hour, from \texttt{p\_flood\_poly} ($s=6.58$\,m; the model depends on $w/s$ only). At $w/s=1$ a flood has happened. Below 0.20 the level is held at 0.20. $^\dagger$\,Earlier than any storm reached the level: the hour is held at the envelope (Table~\ref{tab:envelope}).}\label{tab:grid}\end{table}

Three things to read from Table~\ref{tab:grid}. At a fixed level, $P$ falls as the
storm goes on, after its first hours. Every row rises with the level. And at the severe level $P$ steps to~1:
just below it, late in a storm, $P$ is well below 1, so an error in the severe
threshold that moves $s$ across the water level moves $P$ by the whole of the
remaining gap.

\subsection{Severe-Threshold Error}
The threshold $s$ is per-gauge and enters only through $h$, so a relative error $e$
in $s$ shifts $h$ by about $-e$, and moves $P$ by about $-e\,\partial P/\partial h$.
\texttt{stress.sensitivity.threshold\_sensitivity} reports this, using the model's own
closed-form slope, including along the envelope where the hour is held.

% ===================================================================
\mkmvalidation
% ===================================================================

\subsection{What the v1.0 Benchmark Found}
The Wolfram benchmark of v1.0 (MKM-FPO-001\_wolfram\_benchmark v0.1, 03-Nov-2026) found
the code computed its document's equations exactly (worst relative error
$1.8\times10^{-14}$), but the document's claims did not hold:
\begin{itemize}[nosep]
    \item $P$ fell as the water rose: at hour~24 it passed $\tfrac12$ at both 52\% and
          74\% of severe.
    \item 8 of Table~1's 54 cells and none of Table~2's six were reproduced.
    \item At the severe level $P$ was 0.00 at hour~0 and 0.29 at hour~167, not 0.42
          and 0.97.
    \item The accuracy claim $|P-\text{GBM}|<0.05$ failed: near severe the 95th
          percentile reached 0.26.
    \item The stress route set $P=1$ at severe without the document saying so.
    \item \texttt{p\_flood\_simple} is a ratio, not a probability.
\end{itemize}
Each of those is answered above: the step at severe in Eq.~\eqref{eq:model}, the
constrained fit, the regenerated tables, the measured accuracy, and \S\ref{sec:simple}.

\subsection{Implementation Benchmark}
Version 2.0 is benchmarked in the same way (MKM-FPO-001\_wolfram\_benchmark v0.2):
the code against an independent 50-digit reference written from this document, Wolfram
against the reference, and this document's tables and claims against the code and the
classifiers.

% ===================================================================
\mkmlimitations
% ===================================================================

\subsection{Assumptions}
\begin{enumerate}[nosep]
    \item \textbf{FPO-A1: One kernel for all gauges.} The features normalise each
          gauge by its own severe threshold, and one kernel is fitted to all 11 gauges
          together. Accuracy differs by gauge (Table~\ref{tab:fit}).
    \item \textbf{FPO-A2: Level and hour suffice.} The rate of rise is left out. It is
          the main source of the remaining difference from the classifiers.
    \item \textbf{FPO-A3: The envelope is the support.} The envelope is measured on
          the storm catalogue. A real storm that rises faster than any in the
          catalogue is priced as the fastest one that did.
\end{enumerate}

\subsection{Limitations}
\begin{enumerate}[nosep]
    \item \textbf{FPO-L1: Accuracy near severe.} The difference from the classifiers
          exceeds 0.1 in \fpNearOverLo\% to \fpNearOverHi\% of states near severe.
    \item \textbf{FPO-L2: Not conservative everywhere.} The bias near severe is
          negative at some gauges.
    \item \textbf{FPO-L3: The step at severe.} $P$ jumps to 1 at the severe level, so
          results near it are sensitive to the threshold (\S\ref{sec:simple} and the
          sensitivity section).
    \item \textbf{FPO-L4: Inherits the classifiers.} It approximates MKM-FC-001 and
          carries whatever error the classifiers carry.
\end{enumerate}

% ===================================================================
\mkmgovernance
% ===================================================================

\subsection{Q1 --- Ronseal Test}
Yes, now. v1.0 did not do what its document said; v2.0 is documented from the code
and its tables are generated from it.

\subsection{Q2 --- Accuracy and Stability}
Near-severe mean absolute difference \fpNearMaeLo{}--\fpNearMaeHi{} against the
classifiers on held-out storms. The surface is monotone in the level by construction.

\subsection{Q3 --- Resource Usage}
Negligible: ten coefficients and the standard library.

\subsection{Q4 --- Override Controls}
None in the model. Coefficients and envelope are configuration, changed by refitting.

\subsection{Q5 --- Output Consumption}
Stress scenarios, the portfolio stress screen and the sensitivity diagnostic.

\subsection{Q6 --- Model Chaining}
Yes: it approximates MKM-FC-001 and feeds the stress pipeline (MKM-ST-001).

\subsection{Q7 --- Weaknesses}
Accuracy near severe, the omitted rate of rise, and the step at severe.

\subsection{Q8 --- Monitoring}
Refit and re-measure when the classifiers are retrained or the storm catalogue
changes; re-run the implementation benchmark on any change to the model or its
configuration.

\subsection{Q9 --- Overall Risk}
Acceptable for stress testing, where it gives a consistent monotone surface across
gauges. It should not stand in for a gauge's classifier where that gauge's own
probability is needed.

% ===================================================================
\mkmchangelog
% ===================================================================

\begin{itemize}[nosep]
    \item \textbf{v2.0:} a new model after the v1.0 benchmark: $P=1$ at or above
          severe; a cubic kernel fitted to all 11 gauges' classifiers on real storm
          states, constrained monotone in the level; the level floor; the support
          envelope; the browser served the model instead of a copy;
          \texttt{p\_flood\_simple} documented as a ratio; tables generated from the
          code.
\end{itemize}

\mkmhistorypage{28-Mar-2026 & 1.0 & Initial documentation: quadratic kernel fitted to GAUGE-dc5e88fb & D. Kelly \\
03-Nov-2026 & 2.0 & New model after the Wolfram benchmark of v1.0; refitted to every gauge & David K Kelly \\}

\end{document}
